\chapter{Finite-field constructions and modular obstructions}
\label{construction:chapter}

This chapter supplies an independent finite-field construction, without
importing code from the separately licensed Singer formalization.
Bose--Chowla gives $p$ Sidon points in $[p^2-1]$ for every prime $p$;
this is sufficient for every construction reduction used in this project.
The sharp all-large-$N$ rate in Corollary~\ref{target:largest} additionally
requires a power-saving prime-in-interval theorem.  A construction along
an unbounded sequence of prime-field lengths does not by itself supply
that uniform number-theoretic input.

\begin{lemma}[Shifted products outside a base field]
\label{construction:shifted-products}
\lean{Sidon.BoseChowla.pair_eq_of_shifted_mul_eq}\leanok
Let $F\subset E$ be fields and $\theta\in E\setminus F$.
If $a,b,c,d\in F$ and
\[
 (\theta+a)(\theta+b)=(\theta+c)(\theta+d),
\]
then either $a=c,b=d$ or $a=d,b=c$.
\end{lemma}
\begin{proof}
The inclusion is the algebra map $\iota:F\to E$, which is injective
because it is a unital homomorphism from a field to a nontrivial field.
Set $u=a+b-c-d$ and $v=ab-cd$.  Expanding the displayed equality and
cancelling $\theta^2$ gives
$\theta\iota(u)+\iota(v)=0$.
If $u\ne0$, injectivity gives $\iota(u)\ne0$, so division gives
$\theta=\iota(-v/u)$, contrary to the choice of $\theta$.
Thus $u=0$.  The same identity now gives $\iota(v)=0$, hence $v=0$.
Consequently $a+b=c+d$ and $ab=cd$.
The identity
\[
 (a-c)(a-d)=a^2-a(c+d)+cd=a^2-a(a+b)+ab=0
\]
shows $a=c$ or $a=d$.  In the first case cancellation in the sum equality
gives $b=d$; in the second, commute $c+d$ first and cancel to get $b=c$.
These alternatives include repeated entries and every characteristic,
including characteristic two.
\end{proof}

\begin{lemma}[A multiplicative Sidon family]
\label{construction:unit-family}
\uses{construction:shifted-products}
\lean{Sidon.BoseChowla.shifted_ne_zero,Sidon.BoseChowla.isSidon_shiftedUnitSet,Sidon.BoseChowla.card_shiftedUnitSet}\leanok
For finite $F\subset E$ and $\theta\notin F$, the units
$\{\theta+x:x\in F\}$ form a multiplicative Sidon family of size $|F|$.
\end{lemma}
\begin{proof}
For each $x\in F$, equality $\theta+\iota(x)=0$ would give
$\theta=\iota(-x)$, so this element is nonzero and defines a unit of $E$.
Equality between the units for $x$ and $y$ implies equality between their
underlying field elements; cancelling $\theta$ and using injectivity of
$\iota$ gives $x=y$.  The image of the finite universe of $F$ therefore
has cardinality $|F|$.
Interpret the commutative multiplicative unit group as an additive group
using the type synonym \texttt{Additive}.  Equality of two sums there is
equality of two products of units.  Coercing to $E$ and applying
Lemma~\ref{construction:shifted-products} gives precisely equality of
the corresponding unordered pairs.  This proves the Sidon property of
the entire finite family.
\end{proof}

\begin{lemma}[Cyclic groups to integer intervals]
\label{construction:cyclic-lift}
\lean{IsSidonSet.image_addEquiv,Sidon.isSidon_liftResidues,Sidon.card_liftResidues,Sidon.exists_integer_sidon_of_cyclic_group}\leanok
A finite Sidon subset $A$ of a finite cyclic additive group $G$ lifts to a
Sidon subset of $[|G|]$ with cardinality $|A|$.
\end{lemma}
\begin{proof}
Put $M=|G|$.  A group is nonempty, so $M>0$.
The standard cyclic-group classification gives an additive isomorphism
$e:\mathbb Z/M\mathbb Z\to G$; in Lean this is
\texttt{zmodAddCyclicAddEquiv}.
The image $R=e^{-1}(A)$ has cardinality $|A|$.
An equality of sums in $R$ maps under $e$ to an equality of sums in $A$.
The Sidon property gives equality of the image pairs; injectivity of $e$
gives equality of the original pairs.  Thus $R$ is Sidon.

For each residue $r$, let $r_0\in\{0,\ldots,M-1\}$ be its canonical
integer representative and set
$S=\{r_0+1:r\in R\}$.  The inverse identity
$r_0\bmod M=r$ proves that the representative map is injective, so
$|S|=|R|$.  Its values satisfy $1\le r_0+1\le M$.
If two sums from $S$ are equal, cancel the two added ones and reduce the
resulting integer equality modulo $M$.  This gives equality of two sums
in $R$, so the residue pairs are equal.  Apply the representative-plus-one
map to those pairs to recover equality of the integer pairs.
No bound on the sums themselves is required: reduction modulo $M$ is used
only in the direction from integer equality to residue equality.
\end{proof}

\begin{theorem}[Prime-field Bose--Chowla construction]
\label{construction:bose-chowla}
\uses{construction:unit-family,construction:cyclic-lift}
\lean{Sidon.BoseChowla.exists_shift_outside_base,Sidon.BoseChowla.exists_integer_sidon,Sidon.BoseChowla.exists_prime_sidon}\leanok
For every prime $p$, there is a Sidon set $S\subset[p^2-1]$ with $|S|=p$.
\end{theorem}
\begin{proof}
Take $F=\mathbb F_p=\texttt{ZMod}\ p$ and
$E=\mathbb F_{p^2}=\texttt{GaloisField}\ p\ 2$.
Mathlib constructs the latter as a splitting field and proves it is a
field, an $F$-algebra, and has cardinality $p^2$.
Since $p\ge2$, we have $p<p^2$.
If the algebra map $F\to E$ were surjective, the finite-cardinality
inequality for a surjection would give $p^2\le p$, a contradiction.
Choose $\theta$ outside its range and apply
Lemma~\ref{construction:unit-family}.

The multiplicative group $E^\times$ is cyclic by the finite-field cyclicity
theorem already present in pinned mathlib.  Its cardinality is
$|E|-1=p^2-1$: its elements correspond exactly to the nonzero elements of
$E$.  Lemma~\ref{construction:cyclic-lift}, applied after the additive type
synonym, gives the claimed integer set and exact cardinality.
The proof uses no discrete logarithm computation or decidable search.
\end{proof}

\begin{lemma}[Affine modular obstruction]
\label{construction:affine-obstruction}
\uses{construction:bose-chowla}
\lean{Sidon.card_residueImage,Sidon.isSidon_residueImage,Sidon.residueImage_subset_interval,Sidon.residueImage_no_solution,Sidon.exists_prime_residue_obstruction}\leanok
Let $q\ge1$ and let integer coefficients $a_i$ satisfy
$q\nmid\sum_i a_i$.  If $S_0\subset[L]$ is Sidon, then
$S=qS_0+1\subset[qL+1]$ is Sidon, has size $|S_0|$, and contains no
solution to $\sum_i a_i x_i=0$.
In particular each prime $p$ supplies such a set of cardinality $p$ in
$[q(p^2-1)+1]$.
\end{lemma}
\begin{proof}
The map $y\mapsto qy+1$ is injective because multiplication by the nonzero
integer $q$ is cancellable.  Equality of two sums in its image first
cancels the two added ones and then multiplication by $q$, giving an
equality of two sums from $S_0$.  Its Sidon property proves the image is
Sidon, and injectivity preserves cardinality.
For $1\le y\le L$ we have $1\le qy+1\le qL+1$.
If all $x_i=qy_i+1$, finite distributivity gives
\[
 \sum_i a_i x_i=q\sum_i a_i y_i+\sum_i a_i.
\]
Were the left side zero, this equality would exhibit
$\sum_i a_i=q(-\sum_i a_i y_i)$, contradicting the divisibility hypothesis.
The Bose--Chowla set supplies $L=p^2-1$ and $|S_0|=p$.
\end{proof}

\begin{lemma}[Exact sharp-density subsequences]
\label{construction:subsequence-obstructions}
\uses{construction:affine-obstruction}
\lean{Sidon.exists_prime_odd_obstruction,Sidon.exists_arbitrarily_large_odd_obstruction,Sidon.exists_arbitrarily_large_residue_obstruction}\leanok
There are arbitrarily large lengths $N$ with a Sidon subset of $[N]$ of
size at least $\sqrt{N}/\sqrt2$ containing no solution to
$x_1+x_2+x_3+x_4=x_5$.
For any integer coefficients with nonzero sum, some fixed prime $q$
supplies arbitrarily large lengths with an equation-free Sidon subset of
size at least $\sqrt{N}/\sqrt q$.
\end{lemma}
\begin{proof}
For parity, take $N=2(p^2-1)+1$ and $S=2S_0+1$ from
Theorem~\ref{construction:bose-chowla}.  All elements are odd, whereas a
sum of four odd integers is even, so the equation is impossible even with
repeated variables.  The set has cardinality $p$, and
$N=2p^2-1\le2p^2$ gives $\sqrt N/\sqrt2\le p$.
Primes are unbounded; hence these lengths are unbounded.
The extra $+1$ in $N$ is necessary because the affine map adds one.

For a nonzero coefficient sum $\sigma$, choose a prime $q>|\sigma|$.
Such a prime exists by unboundedness of primes.  It cannot divide
$\sigma$: otherwise $q\mid|\sigma|$ and positivity of $|\sigma|$ would
imply $q\le|\sigma|$.  Apply the affine obstruction with
$N=q(p^2-1)+1$.  Since $q\ge1$, $N\le qp^2$, giving the stated lower
bound.  Varying the independent construction prime $p$ makes $N$
unbounded while the obstruction prime $q$ stays fixed.
These subsequence statements do not claim the sharp coefficient for all
large $N$ in Proposition~\ref{app:odd-obstruction}.
\end{proof}

\begin{theorem}[The positive-density conclusion in paper Remark 1.14]
\label{construction:dense-modular-obstruction}
\uses{construction:affine-obstruction}
\lean{Sidon.exists_prime_construction_below,Sidon.exists_dense_residue_obstruction}\leanok
For any finite family of integer coefficients with nonzero sum, there is
a prime $q$ such that every $N\ge4q+1$ admits an equation-free Sidon set
$S\subset[N]$ with
\[
 |S|\ge\frac{\sqrt N}{2\sqrt q}.
\]
\end{theorem}
\begin{proof}
Choose the obstruction prime $q$ as above.  Put
$M=\lfloor(N-1)/q\rfloor$, $\ell=\lfloor\sqrt M\rfloor$, and
$n=\lfloor\ell/2\rfloor$.
The threshold gives $M\ge4$, hence $\ell\ge2$ and $n\ge1$.
Mathlib's Bertrand theorem supplies a prime $p$ with
$n<p\le2n\le\ell$.
Thus $p^2\le\ell^2\le M$, so
$q(p^2-1)+1\le qM+1\le N$ and the affine Bose--Chowla set embeds in
$[N]$.
For the density bound, $n<p$ gives $\ell+1\le2p$.
The defining floor-square-root inequality gives
$M<(\ell+1)^2\le4p^2$, hence $M+1\le4p^2$.
Euclidean division of $N-1$ by $q$ gives
$N-1=qM+r$ with $0\le r<q$, so $N\le q(M+1)\le4qp^2$.
Taking nonnegative square roots and dividing by $2\sqrt q>0$ yields
$\sqrt N/(2\sqrt q)\le p=|S|$.
This proves the paper's positive-density assertion at every sufficiently
large length using only Bertrand, without a prime-number-theorem input.
\end{proof}


\section{Prime density and sharp all-length obstructions}

\begin{lemma}[The checked analytic prime-density input]
\label{construction:prime-number-theorem}
\lean{Sidon.NumberTheory.Tauberian.WeakPNT,Sidon.NumberTheory.psi_isEquivalent,Sidon.NumberTheory.theta_isEquivalent}\leanok
Let $\Lambda$ denote the von Mangoldt function and put
$\psi(x)=\sum_{n\le x}\Lambda(n)$ and
$\theta(x)=\sum_{p\le x,\ p\ {\rm prime}}\log p$.
Then $\psi(x)\sim x$ and $\theta(x)\sim x$ as $x\to\infty$.
\end{lemma}
\begin{proof}
The local proof ports the proved Wiener--Ikehara argument from the
Apache-licensed PrimeNumberTheoremAnd project, with its unused
bounded-variation alternative removed.  Here are the analytic steps,
including the hypotheses needed at each limit.

For a nonnegative arithmetic function $f$, assume
$\sum_{n<N}f(n)\le CN$, absolute convergence of
$F(s)=\sum_{n\ge1}f(n)n^{-s}$ for $\Re s>1$, and continuous extension of
$G(s)=F(s)-A/(s-1)$ to $\Re s\ge1$.
For a compactly supported $C^2$ function $v$ on $\mathbb R$, two
integrations by parts give
\[
 (1+u^2)\widehat v(u)=
 \widehat{v-v''/(4\pi^2)}(u),\qquad
 |\widehat v(u)|\le
 \frac{\|v\|_1+(4\pi^2)^{-1}\|v''\|_1}{1+u^2}.
\]
Expanding the absolutely convergent Dirichlet series under the integral
gives, for $\sigma>1$,
\[
 \sum_{n\ge1}f(n)n^{-\sigma}
   \widehat v\!\left(\frac{\log(n/x)}{2\pi}\right)
 =\int_{\mathbb R}F(\sigma+it)v(t)x^{it}\,dt.
\]
The sum of absolute values is dominated by
$\sum f(n)n^{-\sigma}\int|v|$, which is finite, so Fubini applies.
Computing $(\sigma+it-1)^{-1}$ as an exponential integral and subtracting
its contribution gives an identity with $G$ on the right.
Summation by parts, the bound $CN$, and the displayed Fourier decay
supply a summable bound uniform as $\sigma\downarrow1$.
On the compact support of $v$, continuity of $G$ gives the corresponding
integrable bound.  Dominated convergence therefore gives
\[
 \sum_{n\ge1}\frac{f(n)}n
   \widehat v\!\left(\frac{\log(n/x)}{2\pi}\right)
 -A\int_{-\log x}^{\infty}\widehat v(u/(2\pi))\,du
 =\int_{\mathbb R}G(1+it)v(t)x^{it}\,dt.
\]
The right side tends to zero by the Riemann--Lebesgue lemma at frequency
$-\log x/(2\pi)$.

Fourier inversion and the logarithmic substitution turn this into
$x^{-1}\sum_n f(n)\Phi(n/x)\to A\int_0^\infty\Phi$ for smooth $\Phi$
compactly supported in $(0,\infty)$.
The required extension of compact Fourier tests is justified by
multiplying a function with integrable first two derivatives by smooth
cutoffs at scale $R$: the product rule and dominated convergence give
convergence in the norm
$\|v\|_1+(4\pi^2)^{-1}\|v''\|_1$.
Smooth upper and lower approximations to $1_{[a,b]}$, with $0<a<b$,
and nonnegativity of $f$ now give interval mass $A(b-a)$.
For $[a,b]=[\varepsilon,1]$ the omitted mass is at most
$C\lceil\varepsilon N\rceil/N$; first take $N\to\infty$ and then
$\varepsilon\downarrow0$.  This proves
$N^{-1}\sum_{n<N}f(n)\to A$.

Apply this theorem to $f=\Lambda$ and $A=1$.
Mathlib supplies $\Lambda\ge0$, its Chebyshev linear bound and absolute
L-series convergence, and the continuous auxiliary function equal to
$-\zeta'(s)/\zeta(s)-1/(s-1)$ on $\Re s>1$.
Its continuity on the boundary uses the checked nonvanishing theorem for
$\zeta$ on $\Re s=1$ away from the pole.
The extra endpoint term in passing from $n<N$ to $n\le N$ is bounded by
$\log N/N\to0$.
Composing with $\lfloor x\rfloor$ and using
$0\le x-\lfloor x\rfloor<1$ proves $\psi(x)\sim x$.
Finally
\[
 0\le\psi(x)-\theta(x)\le2\sqrt x\log x=o(x)
\]
gives $\theta(x)\sim x$.
Its dependency audit reports only the three standard Lean axioms:
\[
 \texttt{propext},\quad \texttt{Classical.choice},\quad \texttt{Quot.sound}.
\]
\end{proof}

\begin{lemma}[Primes in every sufficiently large relative interval]
\label{construction:relative-primes}
\uses{construction:prime-number-theorem}
\lean{Sidon.exists_prime_of_theta_lt,Sidon.eventually_prime_in_relative_interval,Sidon.relativePrimeIntervals}\leanok
For every $\varepsilon>0$, every sufficiently large real $x$ admits a
prime $p$ with $(1-\varepsilon)x\le p\le x$.
\end{lemma}
\begin{proof}
First suppose $0<\varepsilon<1$ and put $a=(1-\varepsilon)x$.
As $x\to\infty$, both $x$ and $a$ tend to infinity.
The preceding asymptotic gives
$|\theta(x)-x|\le(\varepsilon/4)x$ and
$|\theta(a)-a|\le(\varepsilon/4)a$ for all sufficiently large $x$.
Hence
\[
 \theta(x)-\theta(a)\ge
 \left(\varepsilon-\frac{\varepsilon}{4}(2-\varepsilon)\right)x
 =\left(\frac{\varepsilon}{2}+\frac{\varepsilon^2}{4}\right)x>0.
\]
If no prime belonged to $(a,x]$, then the finite prime sets up to
$\lfloor a\rfloor$ and $\lfloor x\rfloor$ would coincide.
Their logarithmic sums would be equal, a contradiction.
For arbitrary $\varepsilon>0$, apply this result to
$\min(\varepsilon,1/2)$; its interval is contained in the required one.
\end{proof}

\begin{theorem}[Sharp all-length modular obstructions]
\label{construction:sharp-obstructions}
\uses{construction:bose-chowla,construction:affine-obstruction,construction:relative-primes}
\lean{Sidon.eventually_exists_prime_construction_below,Sidon.eventually_exists_sharp_odd_obstruction,Sidon.exists_eventually_sharp_residue_obstruction}\leanok
For every $\delta>0$ and every sufficiently large $N$, there is a Sidon
set $S\subset[N]$ with
$|S|\ge(1/\sqrt2-\delta)\sqrt N$ and no solution to
$x_1+x_2+x_3+x_4=x_5$.
More generally, if integer coefficients have nonzero sum, there is a
fixed prime $q$ such that for every $\delta>0$ and every sufficiently
large $N$, the equation has no solution in a Sidon subset of $[N]$ of
size at least $(1/\sqrt q-\delta)\sqrt N$.
\end{theorem}
\begin{proof}
Fix $q\ge1$ and $\delta>0$, and put
$x=\sqrt{N/q}$ and $\varepsilon=\delta\sqrt q$.
Since $x\to\infty$ with $N$, the relative-prime theorem supplies a prime
$p$ with $(1-\delta\sqrt q)\sqrt{N/q}\le p\le\sqrt{N/q}$.
The upper bound gives $qp^2\le N$, and
$q(p^2-1)+1\le qp^2$ because $q\ge1$.
Thus the affine Bose--Chowla construction embeds in $[N]$.
It has exactly $p$ elements, and
\[
 p\ge(1-\delta\sqrt q)\frac{\sqrt N}{\sqrt q}
   =(1/\sqrt q-\delta)\sqrt N.
\]
For Proposition~\ref{app:odd-obstruction}, take $q=2$ and use the parity
obstruction.  For general coefficients choose a prime $q$ not dividing
their nonzero sum, then use the modular obstruction.
The threshold may depend on $\delta$ and $q$, exactly as the asymptotic
quantifiers require.
\end{proof}


\section{The quantitative prime-interval frontier}
\label{construction:prime-frontier}

The relative-prime theorem above suffices for the sharp obstruction
constants, but it does not provide the exponent required by
Corollary~\ref{target:largest}.
The checked reduction below retains the explicit prime-interval premise.
The \href{https://github.com/TomasOrtega/Sidon/blob/master/docs/ingham-prime-interval-proof.md}{expanded analytic derivation}
records the proof being implemented; the following statements give its
checked finite interfaces and exact remaining inputs.

\begin{lemma}[Checked reduction for largest Sidon sets]
\label{construction:largest-reduction}
\uses{construction:bose-chowla,target:uniformity}
\lean{Sidon.exists_maximum_sidon,Sidon.exists_sidon_near_sqrt_of_prime_interval,Sidon.IsMaximumSidon.deviation_le_of_prime_interval,Sidon.IsMaximumSidon.fourier_uniformity_of_prime_interval,Sidon.largest_sidon_fourier_uniformity_of_prime_interval}\leanok
Suppose $C\ge0$ and every $x\ge X$ has a prime
$x-Cx^{2/3}\le p\le x$.
Then the conclusion of Corollary~\ref{target:largest} holds, with
Fourier constant $32\sqrt{C+5}$ for all sufficiently large $N$.
\end{lemma}
\begin{proof}
A maximum-cardinality Sidon subset exists because the finite power set
of $[N]$ has the empty set as an admissible member. For
$\sqrt N\ge X$, choose a prime
$\sqrt N-CN^{1/3}\le p\le\sqrt N$.
The Bose--Chowla set lies in $[p^2-1]\subset[N]$ and has size $p$.
Every maximum Sidon set therefore has size
$m\ge\sqrt N-CN^{1/3}$.
The checked sharp size bound gives $m\le\sqrt N+4N^{1/4}$.
Since $N\ge1$ implies $N^{1/4}\le N^{1/3}$, dividing these inequalities
by the positive $\sqrt N$ gives
\[
 |m/\sqrt N-1|\le(C+4)N^{-1/6}.
\]
Theorem~\ref{target:uniformity}, with its checked constant $16$, gives
\[
 \|\widehat{1_S}-(m/N)\widehat{1_{[N]}}\|_\infty
 \le16\sqrt{C+5}\sqrt N\,N^{-1/12}.
\]
Choose the lower threshold before $N$ so that $\sqrt N\ge X$,
$N\ge1$, and $CN^{-1/6}\le1/2$. The last condition implies
$m\ge\sqrt N/2$. The transform of an indicator has norm exactly $m$:
the triangle inequality bounds every value by $m$, while its value at
zero is $m$. Substitute $\sqrt N\le2m$ in the preceding estimate.
The existence of the prime-interval constants is the remaining
hypothesis; this lemma does not assert it.
\end{proof}

\begin{lemma}[Finite Dirichlet-polynomial mean square]
\label{construction:dirichlet-mean-square}
\lean{Sidon.norm_integral_frequencyChar_le,Sidon.integral_finiteFrequencySum_norm_sq,Sidon.integral_finiteFrequencySum_norm_sq_le}\leanok
For a finite set $I$, complex coefficients $a_i$, distinct real frequencies
$\omega_i$, and $T\ge0$, put $P(t)=\sum_{i\in I}a_i e^{i\omega_i t}$.
Then
\[
 \int_0^T|P(t)|^2\,dt\le T\sum_{i\in I}|a_i|^2+
 \sum_{i\in I}\sum_{\substack{j\in I\\j\ne i}}
       \frac{2|a_i||a_j|}{|\omega_i-\omega_j|}.
\]
\end{lemma}
\begin{proof}
Expand $|P(t)|^2=P(t)\overline{P(t)}$ into the finite double sum.
Continuity makes every summand integrable, so integration commutes with
the sum. The diagonal terms are the constant functions $|a_i|^2$ and
integrate to $T|a_i|^2$. For $i\ne j$,
$\omega_i-\omega_j\ne0$, and the antiderivative
$e^{i(\omega_i-\omega_j)t}/(i(\omega_i-\omega_j))$ gives
\[
 \left|\int_0^T e^{i(\omega_i-\omega_j)t}\,dt\right|
 \le\frac{|e^{i(\omega_i-\omega_j)T}|+1}{|\omega_i-\omega_j|}
 =\frac2{|\omega_i-\omega_j|}.
\]
Multiply by $|a_i||a_j|$ and sum the resulting finite inequalities.
The checked identity precedes the inequality, so no cancellation or
interchange of infinite sums is assumed.
\end{proof}

\begin{definition}[Multiplicity-sensitive finite zero counts]
\label{construction:zero-count}
\lean{Sidon.criticalZetaZeros,Sidon.zetaZeroOccurrences,Sidon.zetaZeroCount,Sidon.analyticOrderAt_riemannZeta_ne_top,Sidon.zeta_zero_multiplicity_pos}\leanok
For real $T$, let $Z_T$ contain one pair $(\rho,k)$ for every actual
zeta zero $\rho=\beta+i\gamma$ with $0\le\beta\le1$,
$0<\gamma\le T$, and $0\le k<m(\rho)$, where $m(\rho)$ is the
analytic order of the zero.  Put
$N(\sigma,T)=|\{(\rho,k)\in Z_T:\sigma\le\beta\}|$.
\end{definition}
\begin{proof}
The relevant zeros lie in the closed complex ball of radius $|T|+2$.
Mathlib proves that the zeta-zero set is closed and discrete, so its
intersection with this compact ball is finite.
The analytic order cannot be infinite: infinite order would make zeta
identically zero on a neighborhood, contradicting discreteness of its
zero set in the nondiscrete complex plane.
At an actual zero, $\rho\ne1$ because the defined value $\zeta(1)$ is
nonzero.  Zeta is analytic there, so its finite analytic order is
positive.  Replace each such zero by its finite range of multiplicity
indices.  This construction requires neither simple zeros nor an
assumed enumeration of them.
\end{proof}

\begin{lemma}[Uniform density controls every finite layer]
\label{construction:zero-layers}
\uses{construction:zero-count}
\lean{Sidon.zeroLayer_card_le_of_density,Sidon.finite_zero_bound_of_density,Sidon.finite_zero_layer_bound,Sidon.zeta_zero_weight_bound}\leanok
Let $L>0$, $\eta,K,D\ge0$, $D<L$, and let $Z$ be a finite multiset
represented by distinct occurrence indices, with real parts
$0\le\beta_z\le1-\eta$.  Suppose, with the same $K,D$ for every
$0\le\sigma\le1$, that
\[
 |\{z\in Z:\sigma\le\beta_z\}|\le K e^{D(1-\sigma)}.
\]
Writing $b=D/L$, one has
\[
 \sum_{z\in Z}e^{-L(1-\beta_z)}
 \le \frac{Ke^b e^{-(1-b)\lfloor L\eta\rfloor}}{1-e^{-(1-b)}}.
\]
\end{lemma}
\begin{proof}
The occurrence $z$ belongs to layer
$j=\lfloor L(1-\beta_z)\rfloor$ and has weight at most $e^{-j}$.
Every occurrence in layer $j$ satisfies
$\beta_z>1-(j+1)/L$.
Apply the density bound at
$\sigma=\max(0,1-(j+1)/L)$: it is at most one, and all layer members
are counted there.  Since $D\ge0$, its exponential factor is at most
$e^{b(j+1)}$, including the layers for which the unclipped threshold
would be negative.
The zero-free inequality implies $j\ge\lfloor L\eta\rfloor$.
Summing the remaining geometric series gives the displayed expression;
its denominator is positive because $b<1$.
For actual zeta zeros take $L=\log x$, $D=A\log T$, and replace $K$
by $K(\log T)^B$ in a uniform density estimate
$N(\sigma,T)\le KT^{A(1-\sigma)}(\log T)^B$.
\end{proof}

\begin{lemma}[A logarithmic reserve absorbs the density loss]
\label{construction:zero-reserve}
\uses{construction:zero-layers}
\lean{Sidon.finite_zero_bound_of_log_reserve}\leanok
Suppose the preceding layer count has $K L^B$ in place of $K$, with
$B\in\mathbb N$, and
$(B+2)\log L\le(1-b)\lfloor L\eta\rfloor$.
Then
\[
 \sum_{z\in Z}e^{-L(1-\beta_z)}
 \le \frac{Ke^b}{1-e^{-(1-b)}}L^{-2}.
\]
\end{lemma}
\begin{proof}
The reserve gives
$e^{-(1-b)\lfloor L\eta\rfloor}\le L^{-(B+2)}$.
Multiply the preceding estimate by $L^B$ and cancel its positive
powers.  Floors therefore cost no uncontrolled asymptotic term.
\end{proof}

\begin{lemma}[Every logarithmic reserve is eventually available]
\label{construction:logarithmic-reserve}
\lean{Sidon.logarithmicZeroFreeWidth,Sidon.tendsto_rpow_div_log_rpow_atTop,Sidon.tendsto_logarithmic_width_reserve,Sidon.eventually_logarithmic_width_reserve}\leanok
For $C>0$ define
\[
 \eta_C(T)=\frac{C}{(\log T)^{2/3}(\log\log T)^{1/3}}.
\]
Then $\eta_C(T)\log T/\log\log T\to\infty$.
For every $\delta>0$ and real $B$, eventually
$B\log\log T\le\delta\eta_C(T)\log T$.
\end{lemma}
\begin{proof}
Put $u=\log T$, so $u\to\infty$. For sufficiently large $T$,
$u>1$ and $\log u>0$, and exponent arithmetic gives
\[
 \frac{\eta_C(T)\log T}{\log\log T}
 =\frac{C u^{1/3}}{(\log u)^{4/3}}.
\]
Mathlib's logarithmic little-oh theorem gives
$(\log u)^{4/3}/u^{1/3}\to0$. This ratio is eventually positive,
so its reciprocal tends to positive infinity; multiplication by $C>0$
preserves that limit. Multiply by $\delta>0$, choose its eventual lower
bound $B$, and multiply by the positive $\log\log T$.
This proves a statement about the displayed width function. Its separate
identification as a zero-free width remains an analytic input until
that theorem is checked.
\end{proof}

\begin{lemma}[Removing prime powers]
\label{construction:prime-powers}
\lean{Sidon.card_nonprimePowers_le,Sidon.nonprimeMangoldtMass_le_real,Sidon.nonprimeMangoldtMass_isLittleO,Sidon.exists_prime_of_mangoldt_increment}\leanok
The nonprime-power contribution to $\psi(N)$ is
\[
 O(\sqrt N\log^2N)=o(N^{2/3}).
\]
If a Mangoldt increment on $(M,N]$ exceeds the checked global
prime-power bound at $N$, then $(M,N]$ contains an actual prime.
\end{lemma}
\begin{proof}
A nonprime prime power $p^k\le N$ has $k\ge2$, $p\le\sqrt N$, and
$k\le\log_2 N$.
Thus there are at most
$\lfloor\sqrt N\rfloor\lfloor\log_2N\rfloor$ such integers, even
if different parameter pairs give duplicates.
Each Mangoldt weight is at most $\log N$.
Two logarithms are negligible against $N^{1/6}$.
If there were no prime in the interval, its theta increment would be
zero.  Its psi increment would then be at most the nonprime-power
mass at $N$, since that mass at $M$ is nonnegative, a contradiction.
\end{proof}

For the remaining analytic assembly, choose $A=23/8$, $T=x^{17/50}$,
and $h=x^{2/3}$.  The numerical margins are
$A(17/50)=391/400<1$ and $x/T=o(h)$.
The required inputs are: a uniform zero-density bound with exponent
$A$; a zero-free width $\eta(T)$ with
$\eta(T)\log T/\log\log T\to\infty$; and a truncated explicit formula
for the psi increment with error
$O(x\log^2x/T+\log x)$ at a height in $[T,2T]$ avoiding zero ordinates.
The zero-layer lemma makes the zero contribution $o(h)$, and the
prime-power lemma then yields an actual prime.
The detailed proof expands the mollifier mean squares, the pole-removing
interpolation, the Jensen/Littlewood zero count, derivative-test
optimization, endpoint conventions, and contour obligations.
The zero-density, zeta-growth, zero-free, and truncated-formula inputs
are not marked checked here; no unconditional power-saving prime-gap
theorem has yet been added to the root dependency graph.

\begin{lemma}[The critical-line mollifier endpoint]
\label{construction:mollifier-endpoint}
\uses{construction:dirichlet-mean-square}
\lean{Sidon.logarithmic_kernel_le,Sidon.sum_range_inv_sqrt_succ_le,Sidon.sum_reciprocal_gap_le,Sidon.integral_criticalMollifier_norm_sq_le,Sidon.integral_dirichletMollifier_critical_norm_sq_le}\leanok
For integers $N\ge1$ and real $T\ge0$, the actual finite mollifier
$M_N(s)=\sum_{1\le n\le N}\mu(n)e^{-s\log n}$ satisfies
\[
 \int_0^T|M_N(1/2+it)|^2\,dt
 \le12(T+N)(1+\log N).
\]
\end{lemma}
\begin{proof}
Use frequencies $-\log n$ and set $a_n=\mu(n)/\sqrt n$
in the finite mean-square lemma.
These frequencies are distinct because the logarithm is injective on
positive reals, and $|a_n|\le1/\sqrt n$ because $|\mu(n)|\le1$.
The diagonal contribution is at most $T H_N$, where
$H_N=\sum_{n=1}^N1/n$.
For $0<m<n$, the elementary inequality
$\log(n/m)\ge1-m/n$ gives
\[
 \frac1{\sqrt m\sqrt n\log(n/m)}
 \le\frac1{\sqrt m\sqrt n}+\frac1{n-m}.
\]
Indeed multiply the positive denominators and use
$m\le\sqrt m\sqrt n$.
The symmetric absolute-value version handles both orderings.
For fixed $m$, split the other integers below and above $m$.
Each distance map $n\mapsto m-n$ or $n\mapsto n-m$ is injective into
$\{1,\ldots,N\}$, so the sum of reciprocal gaps is at most $2H_N$.
Also
\[
 \frac1{\sqrt{n+1}}\le2(\sqrt{n+1}-\sqrt n)
 \quad\Longrightarrow\quad
 \sum_{n=1}^N\frac1{\sqrt n}\le2\sqrt N.
\]
The first inequality follows by multiplying by $\sqrt{n+1}>0$ and
using $(\sqrt{n+1}-\sqrt n)^2\ge0$; the sum telescopes.
Thus the product-kernel double sum is at most $4N$, and the reciprocal
gap double sum is at most $2NH_N$.
The finite mean square is consequently bounded by
$TH_N+8N+4NH_N\le12(T+N)(1+\log N)$, using
$1\le H_N\le1+\log N$.
Finally the exact identity
$e^{-(1/2+it)\log n}=n^{-1/2}e^{-it\log n}$ identifies this finite
frequency sum with the complex-variable mollifier.
\end{proof}

\begin{lemma}[Finite mollifier coefficients]
\label{construction:mollifier-coefficients}
\lean{Sidon.differentiable_dirichletMollifier,Sidon.mollifierCoefficient_eq_one_apply,Sidon.mollifierCoefficient_eq_zero,Sidon.abs_mollifierCoefficient_le}\leanok
The function $M_N$ above is entire.
For the finite divisor coefficient
$a_N(n)=\sum_{d\mid n,\ d\le N}\mu(d)$, one has
$|a_N(n)|\le d(n)$ and, when $n\le N$,
$a_N(n)=1$ for $n=1$ and $a_N(n)=0$ otherwise.
\end{lemma}
\begin{proof}
Each summand of $M_N$ is an exponential of a complex linear function,
so its finite sum is entire.
If $n\le N$, every divisor of $n$ is at most $N$; the coefficient is
therefore the full M\"obius divisor sum, equal to the unit arithmetic
function.  This includes $n=0$, whose divisor finset is empty.
The triangle inequality and $|\mu(d)|\le1$ bound the absolute value by
the number of included divisors, at most $d(n)$.
The infinite Dirichlet-series product with zeta will separately require
absolute convergence on $\Re s>1$; this finite coefficient identity
does not assume an unjustified infinite rearrangement.
\end{proof}
